\documentclass[portrait,a0paper,fontscale=0.30]{baposter}

% Compile with double 'pdflatex poster'

\usepackage[utf8]{inputenc}

% Choose your favorite font here. I go with Fira by default.
%
\usepackage[sfdefault]{FiraSans}          % very recommended
%\usepackage[sfdefault]{cabin}             % pretty good also with 'condensed'
%\usepackage{cmbright}                     % also recommendable, more TeX-ish
%\usepackage[defaultsans]{droidsans}
%\usepackage[condensed,sfdefault]{roboto}  % also works with 'condensed'
%
% The following fonts contain serifs, but are still pretty good for posters.
% Comment out the familydefault setting below if using serif fonts.
%
%\usepackage[mono=false]{libertine}  % Libertine is almost standard now.
%\usepackage{fouriernc}              % URW schoolbook, for lovable vintage-style posters.
%\usepackage{lmodern}                % The usual TeX font for TeX and math-oriented posters.

\usepackage[T1]{fontenc}
\usepackage{color}

% Comment this out if using serif fonts.
\renewcommand{\familydefault}{\sfdefault}

% Increase/decrease this as you like; the units are in the baposter-ish scale.
\newcommand{\elixirFooterSize}{4cm}

\definecolor{eblue}{RGB}{0,81,113}
\definecolor{eorange}{RGB}{238,117,2}
\color{eblue}

\begin{document}

% Custom poster background; extra coordinate is set for other nodes to follow below.
\background{\begin{tikzpicture}[remember picture,overlay]
    \node [inner sep=0pt, anchor=north west] (headerbg)
      at (current page.north west)
      {\includegraphics[width=\paperwidth]{header.jpg}};
    \coordinate (headerbgse)
      at ([yshift=-1em] headerbg.south east);
    \coordinate (footerstart)
      at ([yshift=\elixirFooterSize] current page.south);
    \fill [fill=eorange]
      (current page.west |- footerstart)
      rectangle
      ([yshift=-1em] current page.east |- footerstart);
    \coordinate (footersepse)
      at ([yshift=-2em] footerstart);

    % small notice about where to get this template; omit if needed.
    \node [anchor=south west, rotate=90, text width=10em, font=\tiny, color=black!25]
      at ([xshift=-1pt, yshift=1pt] current page.south east)
      {\TeX\ ELIXIR poster template \\ \texttt{github.com/exaexa/elixirposter-tex}};
\end{tikzpicture}}

\begin{poster}{%
  eyecatcher=false,
  background=user,
  columns=2,
  headerborder=none,
  textborder=none,
  headershape=rectangle,
  headershade=plain,
  headerColorOne=white,
  boxshade=none,
  headerFontColor=eorange,
  headerfont=\Large
  }{}%
  {\vspace{2em} \color{white} \fontsize{32pt}{32pt}\selectfont
  Put the poster title here. %TODO
  }%
  {
  % You may also include a subtitle (not recommended in ELIXIR template)
  }{}


%
% TODO
%
% Poster editing starts here!
%

\begin{posterbox}[name=abstract, column=0, span=2, below=headerbg]{}
Put the long abstract here; template recommends that it spans 2-3 lines of the whole poster text width. Slightly different designs are possible (e.g.~set the span of this box to 1, and continue the right column below `headerbg').
\end{posterbox}

\begin{posterbox}[name=boxleft, below=abstract, column=0]{Box in the left part of the poster}
Stuff from the left goes here.

Maybe also pictures.
\end{posterbox}
\begin{posterbox}[name=boxright, below=abstract, column=1]{Another box, now more on the right}
Stuff from the right goes here.

Maybe more pictures.
\end{posterbox}

%
% TODO
%
% Footer editing starts here; you have a bit of space to put posterboxes below
% 'footersep'; which is the position of the thick orange line below.
%

\begin{posterbox}[name=footer, below=footersep, column=0]{Contact}
Your Name Here % if including more names, \underline the presenting person

\texttt{your.email@institute.here}

Your Institute, Your Town, Your Country
\end{posterbox}

\begin{posterbox}[name=funding, below=footersep, column=1]{Funding}
Put all your institute logos here.

These should also include EU flag, and the logo of your corresponding ELIXIR node and/or funding source.
\end{posterbox}

\end{poster}
\end{document}
